% !TEX program = xelatex
% ============================================================
%  订单流微观分类与限价订单簿动力学
%  实习汇报 Beamer Slides
%  Compile with: xelatex main.tex (twice for TOC)
% ============================================================

\documentclass[10pt,aspectratio=169]{beamer}

% ---------- Theme ----------
% 使用传统学术汇报主题 Warsaw(经典蓝色,适合研究报告)
\usetheme{Warsaw}
\usecolortheme{default}
\useinnertheme{rectangles}
\useoutertheme{infolines}
\setbeamertemplate{navigation symbols}{}  % 移除底部导航符号

% ---------- Chinese Support ----------
\usepackage{ctex}
\setCJKmainfont{SimSun}     % 宋体（如缺失可改为 Microsoft YaHei / Noto Serif CJK SC）
\setCJKsansfont{SimHei}     % 黑体
\setCJKmonofont{FangSong}   % 仿宋

% ---------- Math & Symbols ----------
\usepackage{amsmath, amssymb, amsfonts}
\usepackage{mathtools}
\usepackage{bm}

% ---------- Tables & Lists ----------
\usepackage{booktabs}
\usepackage{array}
\usepackage{multirow}
\usepackage{multicol}

% ---------- Graphics & Diagrams ----------
\usepackage{graphicx}
\usepackage{tikz}
\usetikzlibrary{arrows.meta, positioning, shapes.geometric, decorations.pathreplacing}

% ---------- Hyperlinks ----------
\usepackage{hyperref}
\hypersetup{colorlinks=true, linkcolor=., urlcolor=blue}

% ---------- Custom Colors ----------
\definecolor{quantred}{RGB}{198,40,40}
\definecolor{quantblue}{RGB}{25,118,210}
\definecolor{quantgreen}{RGB}{46,125,50}

% ---------- Title Page Info ----------
\title{订单流微观分类研究报告}
\subtitle{限价订单簿动力学与高频 Alpha 因子}
\date{2026 年 4 月 9 日}
\author{王超珑}
\institute{JinYi Capital}

% ============================================================
\begin{document}

\maketitle

% --------------- Outline ---------------
\begin{frame}{汇报大纲}
  \tableofcontents[hideallsubsections]
\end{frame}

% ============================================================
\section{引言与研究背景}
% ============================================================

% \begin{frame}{为什么是微观结构?——一场降维打击}

%   \textbf{数据颗粒度的演进:}

%   \begin{center}
%   \begin{tikzpicture}[node distance=1.4cm, every node/.style={font=\small}]
%     \node[draw, rounded corners, fill=blue!10] (a) {日线 / 分钟线};
%     \node[draw, rounded corners, fill=blue!20, right=of a] (b) {Tick (逐笔成交)};
%     \node[draw, rounded corners, fill=blue!30, right=of b] (c) {L2 快照};
%     \node[draw, rounded corners, fill=blue!50, text=white, right=of c] (d) {MBO (逐笔委托)};
%     \draw[->,thick] (a) -- (b);
%     \draw[->,thick] (b) -- (c);
%     \draw[->,thick] (c) -- (d);
%   \end{tikzpicture}
%   \end{center}

%   \vfill

%   \textbf{现代量化的视角转换:}

%   \begin{itemize}
%     \item 市场不再是一个抽象的宏观定价机制
%     \item 而是\alert{由海量异质性订单构成的复杂动力学系统}
%     \item 每一笔订单都携带\textbf{方向、规模}的物理属性,以及\textbf{博弈意图}
%     \item 纳秒级时间戳使我们得以观察金融市场的"基本粒子"
%   \end{itemize}
% \end{frame}

\begin{frame}{研究的两大核心目标}

  对海量订单流进行微观分类与解构,服务于两个目的:

  \vfill

  \begin{columns}[t]
    \column{0.48\textwidth}
    \begin{block}{① 挖掘 Alpha 因子}
      \begin{itemize}
        \item 构建预测短期价格变动的高频信号
        \item 识别市场中的结构性失衡
        \item 提取不同参与者的行为模式
        \item 优化算法执行降低冲击成本
      \end{itemize}
    \end{block}

    \column{0.48\textwidth}
    \begin{block}{② 防范微观系统性风险}
      \begin{itemize}
        \item 监测有毒订单流的累积
        \item 预警流动性枯竭与闪电崩盘
        \item 识别价格操纵行为
        \item 实时风控做市商敞口
      \end{itemize}
    \end{block}
  \end{columns}

  \vfill

  \begin{alertblock}{核心信念}
    显微镜可以穿透价格变动的表象,揭示\textbf{参与者动机、信息不对称、流动性失衡}三个关键维度。
  \end{alertblock}

\end{frame}

\begin{frame}{五维度分类全景图}
  \begin{center}
  \begin{tikzpicture}[node distance=1.0cm and 0.6cm, font=\scriptsize]
    \node[draw, rounded corners, fill=quantblue!20, minimum width=3cm] (root) {一笔订单事件};

    \node[draw, rounded corners, fill=blue!10, below left=1.0cm and 2.5cm of root] (d1) {① 方向推定};
    \node[draw, rounded corners, fill=blue!10, below left=1.0cm and 0.5cm of root] (d2) {② 流动性力学};
    \node[draw, rounded corners, fill=blue!10, below=1.0cm of root] (d3) {③ 信息毒性};
    \node[draw, rounded corners, fill=blue!10, below right=1.0cm and 0.5cm of root] (d4) {④ 参与者画像};
    \node[draw, rounded corners, fill=blue!10, below right=1.0cm and 2.5cm of root] (d5) {⑤ 失衡因子};

    \draw[->] (root) -- (d1);
    \draw[->] (root) -- (d2);
    \draw[->] (root) -- (d3);
    \draw[->] (root) -- (d4);
    \draw[->] (root) -- (d5);

    \node[below=0.2cm of d1, text width=2.2cm, align=center] {Lee-Ready\\BVC};
    \node[below=0.2cm of d2, text width=2.2cm, align=center] {Maker/Taker\\Kyle 模型};
    \node[below=0.2cm of d3, text width=2.2cm, align=center] {PIN / VPIN\\体积时钟};
    \node[below=0.2cm of d4, text width=2.2cm, align=center] {冰山 / 扫单\\Boehmer};
    \node[below=0.2cm of d5, text width=2.2cm, align=center] {OFI / VOI\\DeepLOB};
  \end{tikzpicture}
  \end{center}

  \vfill
  \centering\small
  \emph{接下来的章节将沿着这条主线逐层展开}
\end{frame}

% ============================================================
\section{订单流方向推定}
% ============================================================

\begin{frame}{方向推定:谁是这笔成交的发起方?}

  \begin{exampleblock}{一笔成交事件的"双重叙事"}
    \begin{itemize}
      \item 同一笔成交,可以是 A 主动找 B,也可以是 B 主动找 A
      \item \textbf{同一笔成交,在两种解释下产生的微观信号完全相反}
      \item \alert{没有方向 → 没有失衡 → 没有 Alpha}
    \end{itemize}
  \end{exampleblock}

  \vfill

  我们必须从价格、时间、报价等周边线索\textbf{反推出真实发起方}。

  这是构建所有后续微观结构因子的\textbf{先决条件}。

\end{frame}

\begin{frame}{Lee-Ready 算法 (1991) ——黄金标准\footnotemark}

  长期占据交易方向推断的\textbf{学术与工业黄金标准}。结合两层规则:

  \vfill

  \begin{block}{第一层:报价测试 (Quote Rule)}
    \begin{itemize}
      \item $P_t > M_t$:成交价高于中间价 → \textcolor{quantgreen}{\textbf{主动买入}}
        \begin{itemize}\item[] \scriptsize 发起方愿意支付溢价立即成交\end{itemize}
      \item $P_t < M_t$:成交价低于中间价 → \textcolor{quantred}{\textbf{主动卖出}}
        \begin{itemize}\item[] \scriptsize 发起方愿意折价立即出货\end{itemize}
      \item $P_t = M_t$:进入第二层测试
    \end{itemize}
  \end{block}

  \begin{block}{第二层:剔除测试 (Tick Test)}
    \begin{itemize}
      \item Uptick($P_t > P_{t-1}$):主动买入
      \item Downtick($P_t < P_{t-1}$):主动卖出
    \end{itemize}
  \end{block}

  \vfill

  \centering\small\emph{德国市场实测准确率约 72.8\%。但在 HFT 时代,纳秒级时间戳错位会引发系统性偏差。}

  \footnotetext[1]{Lee, C. M. C.; Ready, M. J. \emph{Inferring Trade Direction from Intraday Data}. The Journal of Finance, 1991, 46(2), 733--746. \url{https://doi.org/10.1111/j.1540-6261.1991.tb02683.x}}
\end{frame}

\begin{frame}{BVC 算法——批量订单分类的概率革命}

  Easley、López de Prado、O'Hara (ELO) 提出的\alert{范式转换}级方案 Bulk Volume Classification。

  \begin{block}{核心思想}
    \textbf{摒弃逐笔二元判定}(0 或 1),用\textbf{概率统计}将一段时间内的总成交量\textbf{按比例切分}为主动买入与主动卖出。
  \end{block}

  \vfill

  \[
    V_\tau^B = \sum_{i = \text{last}(\tau-1)+1}^{\text{last}(\tau)} V_i \cdot Z\!\left(\frac{\Delta P_i}{\sigma_{\Delta P}}\right)
  \]

  \begin{itemize}
    \item $Z(\cdot)$:正态分布(或 t 分布)\textbf{累积分布函数 CDF}
    \item 价格平稳 → $Z(0) = 0.5$ → 50/50 切分
    \item 极端拉升 → $Z \to 1$ → ~100\% 主动买
    \item 极端砸盘 → $Z \to 0$ → ~100\% 主动卖
  \end{itemize}

\end{frame}

\begin{frame}{BVC 公式逐项拆解}

  \[
    V_\tau^B = \sum_{i = \text{last}(\tau-1)+1}^{\text{last}(\tau)} V_i \cdot Z\!\left(\frac{\Delta P_i}{\sigma_{\Delta P}}\right)
  \]

  \vfill

  \small
  \begin{tabular}{c|p{9.5cm}}
    \toprule
    \textbf{符号} & \textbf{含义} \\
    \midrule
    $V_\tau^B$ & 第 $\tau$ 个聚合区间(Bucket)内,被分类为\textbf{主动买入}的总交易量 \\
    $\tau$ & 当前进行方向估计的批量区间索引(一个"桶") \\
    $i$ & 桶内各子时间柱(Bar)的迭代索引,从上一桶末尾的下一个开始 \\
    $V_i$ & 第 $i$ 个 Bar 的\textbf{总物理交易量}(买+卖的全部) \\
    $\Delta P_i$ & 第 $i$ 个 Bar 内的价格变动(收盘价 $-$ 开盘价) \\
    $\sigma_{\Delta P}$ & 价格变动的样本标准差(短期市场波动率的衡量基准) \\
    $Z(\cdot)$ & \textbf{标准正态分布的 CDF}:将标准化得分映射到 $[0,1]$ \\
    \bottomrule
  \end{tabular}

  \vfill

  \begin{block}{公式运作的直觉}
    \begin{enumerate}
      \item 将价格变动 $\Delta P_i$ 除以波动率 $\sigma_{\Delta P}$ → 得到"\textbf{标准化的价格异动强度}"
      \item 用 CDF 将该强度映射到 $[0,1]$ → 得到"\textbf{这个 Bar 中属于主动买入的比例}"
      \item 用该比例乘以 Bar 总量 $V_i$ → 得到"\textbf{该 Bar 中的主动买入量}"
      \item 对桶内所有 Bar 求和 → 得到"\textbf{整个桶的主动买入量}" $V_\tau^B$
    \end{enumerate}
  \end{block}

  \vfill

  \centering\footnotesize
  主动卖出量自然为 $V_\tau^S = V - V_\tau^B$(守恒定律:桶内总量固定为 $V$)

\end{frame}

\begin{frame}{Lee-Ready vs BVC:速度与精度的权衡}

  \small
  \begin{tabular}{p{3.2cm}|p{4.0cm}|p{4.5cm}}
    \toprule
    \textbf{特征} & \textbf{Lee-Ready / Tick Rule} & \textbf{BVC 批量订单分类} \\
    \midrule
    数据颗粒度 & 逐笔成交 + 纳秒级报价 & K 线 (Bar Data) 即可 \\
    分类哲学 & 二元离散判定 & 连续概率切分 \\
    计算复杂度 & 极高 & 极低 \\
    高频适应性 & 易受时间戳错位干扰 & 天然平滑微观噪音 \\
    微观准确度 & 较高 & 略低 \\
    主要场景 & 学术研究、超高频做市 & 实时风控、闪电崩盘预警 \\
    \bottomrule
  \end{tabular}

  \vfill

  \begin{alertblock}{实战建议}
    \begin{itemize}
      \item 学术研究 / 高精度因子验证 → \textbf{Lee-Ready}
      \item 缺乏 Tick 数据的中频策略 → \textbf{BVC}
      \item VPIN 与闪电崩盘预警 → \textbf{BVC}(它就是 VPIN 的方向分类底层)
    \end{itemize}
  \end{alertblock}

\end{frame}

% ============================================================
\section{流动性力学与价格冲击}
% ============================================================

\begin{frame}{Maker 与 Taker:微观参与者的两种角色}

  在 LOB 主导的现代交易所,所有参与者被映射为两种物理角色:

  \vfill

  \begin{columns}[t]
    \column{0.48\textwidth}
    \begin{block}{Maker(流动性提供者)}
      \begin{itemize}
        \item 提交\textbf{限价单 (Limit Order)}
        \item 价格未与盘口形成交叉
        \item 静静挂在 LOB 等待成交
        \item 构筑"盘口厚度"与防御纵深
        \item \textcolor{quantgreen}{交易所返佣 (Rebate)}
      \end{itemize}
    \end{block}

    \column{0.48\textwidth}
    \begin{block}{Taker(流动性消耗者)}
      \begin{itemize}
        \item 提交\textbf{市价单 (Market Order)}
        \item 或激进的穿透型限价单
        \item 即时与 Maker 撮合成交
        \item 瞬间抽离系统流动性
        \item \textcolor{quantred}{支付手续费 (Fee)}
      \end{itemize}
    \end{block}
  \end{columns}

  \vfill

  \centering\small
  这种角色分类是各大交易所 \textbf{Maker-Taker Fee Model} 的商业基础。

\end{frame}

\begin{frame}{Kyle (1985) 模型:三方博弈的理论基石}

  \textbf{三方参与者:}

  \begin{itemize}
    \item \textcolor{quantred}{\textbf{知情交易者}}:无成本观察真实价值 $v \sim N(p_0, \Sigma_0)$
    \item \textcolor{quantblue}{\textbf{噪音交易者}}:总订单流 $u \sim N(0, \sigma_u^2)$
    \item \textcolor{quantgreen}{\textbf{做市商}}:只能观察 $y = x + u$,设定竞争性出清价格
  \end{itemize}

  \vfill

  \begin{block}{做市商的线性定价规则}
    \[ p = \mu + \lambda y \]
  \end{block}

  \begin{block}{知情交易者的最优策略 (一阶条件求解)}
    \[ x = \frac{v - \mu}{2\lambda} \]
  \end{block}

  \vfill

  \centering\small\emph{知情者克制自己的进攻——价格冲击越大,他越要分散投放}

\end{frame}

\begin{frame}{Kyle's Lambda——市场深度的数学定义}

  贝叶斯均衡解:

  \[
    \boxed{ \lambda = \frac{1}{2} \sqrt{\frac{\Sigma_0}{\sigma_u^2}} }
  \]

  \vfill

  \begin{columns}[t]
    \column{0.48\textwidth}
    \begin{block}{波动 ↑ → 流动性 ↓}
      $\lambda \propto \sqrt{\Sigma_0}$
      \begin{itemize}
        \item 资产先验波动越大
        \item → 做市商逆向选择风险越高
        \item → 加宽价差以自我保护
        \item → 流动性枯竭
      \end{itemize}
    \end{block}

    \column{0.48\textwidth}
    \begin{block}{噪音 ↑ → 流动性 ↑}
      $\lambda \propto 1/\sqrt{\sigma_u^2}$
      \begin{itemize}
        \item 噪音交易越多
        \item → 知情者获得"完美伪装"
        \item → 做市商感知风险被稀释
        \item → 流动性极佳
      \end{itemize}
    \end{block}
  \end{columns}

  \vfill

  \begin{alertblock}{核心洞察}
    任何 Taker 在消耗流动性时,本质上都在向 Maker 支付由 $\lambda$ 决定的\textbf{信息风险溢价}。
  \end{alertblock}

\end{frame}

% ============================================================
\section{信息毒性与高频风控}
% ============================================================

\begin{frame}{有毒订单流 (Toxic Order Flow)}

  \begin{exampleblock}{做市商的痛点}
    \textbf{最忌惮的不是单纯的价格波动}(波动是中性的、可保的),而是不幸与那些\textbf{掌握私有信息、有坚定方向性判断}的资金互为对手盘。
  \end{exampleblock}

  \vfill

  这类"有毒订单流"具有三大特征:

  \begin{itemize}
    \item \textbf{极强的持续性}——信息一旦释放,单向流持续多个时间窗口
    \item \textbf{强大的穿透力}——不在乎成本,只为了建仓
    \item \textbf{方向的确定性}——必然朝一个方向走,做市商必然亏损
  \end{itemize}

  \vfill

  \centering\small
  类比:做市商就像保险公司,而你卖给他房屋险时\textbf{你家已经着火了}。

\end{frame}

\begin{frame}{PIN 模型 (1996) ——开山之作}

  Easley、Kiefer、O'Hara、Paperman 创立的 \textbf{Probability of Informed Trading}。

  \vfill

  \[
    \text{PIN} = \frac{\alpha \mu}{\alpha \mu + \epsilon_b + \epsilon_s}
  \]

  \begin{itemize}
    \item $\alpha$:发生信息事件的先验概率
    \item $\mu$:知情交易者的泊松到达率
    \item $\epsilon_b, \epsilon_s$:不知情噪音买/卖单的基准到达率
  \end{itemize}

  \vfill

  \begin{alertblock}{工业落地的两大瓶颈}
    \begin{itemize}
      \item 参数估计需要 \textbf{MLE 迭代},巨量样本下经常\textbf{不收敛}
      \item 建立在\textbf{时间时钟}上,面对毫秒级毒性流\textbf{迟钝且充满噪音}
    \end{itemize}
  \end{alertblock}

\end{frame}

\begin{frame}{VPIN (2012) ——体积时钟创新}

  Easley、López de Prado、O'Hara 的里程碑式创新:\textbf{彻底抛弃物理时间}。

  \begin{block}{Volume Clock 的物理直觉}
    以"\textbf{每成交固定数量}"而非"每过固定时间"作为采样单位。
    \begin{itemize}
      \item 开盘 30 分钟 → \textbf{极速翻页}(高密度)
      \item 午盘清淡 → \textbf{几乎停滞}
      \item 突发新闻 → \textbf{瞬间翻过多个体积桶}
    \end{itemize}
  \end{block}

  \vfill

  \[
    \text{VPIN}_L = \frac{\sum_{\tau = L-n+1}^{L} |V_\tau^S - V_\tau^B|}{n V}
  \]

  \vfill

  \centering\small
  \textbf{物理类比}:测量心脏健康用"每跳一次心"而非"每秒",VPIN 就是金融市场的心跳计数。
\end{frame}

\begin{frame}{2010 闪电崩盘}

  \textbf{事件}:2010 年 5 月 6 日,道琼斯指数在 5 分钟内自由落体 1000 点,蓝筹股瞬间跌至 1 美分。

  \vfill

  \begin{exampleblock}{VPIN 的预警表现}
    \begin{itemize}
      \item 崩盘前\textbf{一个多小时},VPIN CDF 已突破 \textbf{0.90} 危险阈值
      \item VIX 指数在崩盘前\textbf{毫无异动}——只有 VPIN 提前给出红色警报
    \end{itemize}
  \end{exampleblock}

  \vfill

  \textbf{微观物理过程:}

  \begin{center}
  \footnotesize
  极高 VPIN $\Rightarrow$ 单向有毒流涌入 $\Rightarrow$ 做市商风控触发 $\Rightarrow$
  集体撤单停止做市 $\Rightarrow$ 流动性瞬间抽离 $\Rightarrow$ 少量市价单击穿数百档位 $\Rightarrow$ \alert{系统性崩盘}
  \end{center}

  \vfill

  \begin{alertblock}{结论}
    VPIN 不只是分类指标,更是\textbf{预测极端波动与流动性黑洞的微观地震仪}。
  \end{alertblock}

\end{frame}

% ============================================================
\section{微观资金行为反向工程}
% ============================================================

\begin{frame}{冰山订单 (Iceberg Orders)——机构的隐藏武器}

  \textbf{为什么需要拆单?}大额市价单会因穿透 LOB 多档位产生\alert{毁灭性的市场冲击成本}。

  \vfill

  \begin{block}{冰山订单的物理签名}
    LOB 上只显示极小的 \textbf{Peak Size}(冰山一角),绝大部分隐藏在水面之下(\textbf{Hidden Size})。Peak 被消耗后\textbf{自动 Refill}。
  \end{block}

  \vfill

  \textbf{三种检测算法:}

  \begin{itemize}
    \item \textbf{MBO 数据追踪}:某价位累计成交 $>$ 可见挂单 + 价位未变 → 冰山补充铁证
    \item \textbf{时间差攻击}(针对 Synthetic):微秒级窗口内,前后相同价位、相同量的"幽灵挂单"
    \item \textbf{生存分析}:用 Kaplan-Meier 估计器预测水下隐藏总量
  \end{itemize}

  \vfill

  \centering\small
  \emph{冰山检测是高频量化机构获取短期高胜率 Alpha 的核心利润来源之一。}

\end{frame}

\begin{frame}{扫单流 (Sweep Orders / ISO)——流动性的推土机}

  与冰山的"温吞水"截然相反——\textbf{狂暴、瞬时、推土机式}的订单类型。

  \vfill

  \begin{columns}[T]
    \column{0.5\textwidth}
    \textbf{触发场景:}
    \begin{itemize}
      \item 跨市场套利信号
      \item 宏观经济数据毫秒级释出
      \item ETF 与成分股套利窗口
      \item 期货 vs 现货基差异常
    \end{itemize}

    \column{0.5\textwidth}
    \textbf{行为特征:}
    \begin{itemize}
      \item 不追求价格,只求\textbf{成交率}
      \item 一次性吃掉多个深层档位
      \item 像"流动性真空引爆器"
      \item 触发级联止损的反馈循环
    \end{itemize}
  \end{columns}

  \vfill

  \begin{alertblock}{Alpha 价值}
    识别扫单流 = 微观市场已形成\textbf{不可阻挡的单向惯性}——是高频统计套利的\textbf{第一动量信号}。
  \end{alertblock}

\end{frame}

\begin{frame}{Boehmer 算法 (2021) ——散户订单识别}

  \textbf{颠覆传统假设}:散户不再是噪音,而是新的 Alpha 源泉。

  \vfill

  \begin{block}{底层逻辑:子便士价格改善}
    美国市场的 PFOF(订单流支付)机制下,散户单走\textbf{批发商内部撮合},获得\textbf{零点几美分的价格改善}。机构单则只能在\textbf{整数便士}成交。
  \end{block}

  \vfill

  \small
  \begin{tabular}{p{3.6cm}|p{2.6cm}|p{6cm}}
    \toprule
    \textbf{成交价} & \textbf{推定类型} & \textbf{物理机理} \\
    \midrule
    略低于整数便士 (\$10.999) & 散户主动买 & 试图以 \$11.00 买入,获 \$0.001 折扣 \\
    略高于整数便士 (\$10.001) & 散户主动卖 & 试图以 \$10.00 卖出,获 \$0.001 溢价 \\
    半便士 (\$10.005) & 机构暗池\textbf{(剔除)} & 暗池中点交叉撮合,污染样本 \\
    \bottomrule
  \end{tabular}

  \vfill

  \centering\small\emph{实证发现:散户净买入 $\Rightarrow$ 未来一周约 \textbf{10 bps} 正向 Alpha}
\end{frame}

% ============================================================
\section{基于微观失衡的 Alpha 因子}
% ============================================================

\begin{frame}{OFI——动态事件驱动模型 (Cont et al. 2014)}

  \textbf{颠覆性创新}:不只看成交数据,还看 \textbf{LOB 上的潜在能量}——
  新挂单 + 撤单 + 实际成交\textbf{统一}进事件方程。

  \vfill

  单个事件 $n$ 对买方上行压力的净贡献度:

  \[
  \footnotesize
  e_n = I_{\{P_n^B \geq P_{n-1}^B\}} q_n^B - I_{\{P_n^B \leq P_{n-1}^B\}} q_{n-1}^B
       - I_{\{P_n^A \leq P_{n-1}^A\}} q_n^A + I_{\{P_n^A \geq P_{n-1}^A\}} q_{n-1}^A
  \]

  \vfill

  \begin{itemize}
    \item 第①项:买方提价 / 增量挂单 → \textcolor{quantgreen}{看涨压力 +}
    \item 第②项:买方撤单 / 被砸盘 → \textcolor{quantred}{多方崩溃 -}
    \item 第③项:空方压价 / 囤单 → \textcolor{quantred}{空方压制 -}
    \item 第④项:空方撤单 / 卖价上移 → \textcolor{quantgreen}{防线退缩 +}
  \end{itemize}

  \vfill

  \[ \text{OFI}_k = \sum_n e_n \quad\text{(高频窗口聚合)} \]

\end{frame}

\begin{frame}{OFI 与价格变动:一个堪比物理定律的关系}

  实证发现,在常规微观环境下:

  \[
    \boxed{\Delta P = \beta \cdot \text{OFI} + \epsilon}
  \]

  $R^2$ 在很多个股的微观时间尺度上达到 \textbf{0.6 - 0.8}——远超任何传统技术或基本面因子。

  \vfill

  \begin{block}{极端情况下的非线性缓冲}
    单边失衡过大时,OFI 与 $\Delta P$ 的关系从线性变为\textbf{次线性}:
    \[ |\Delta P| \propto |\Delta V|^\nu, \quad \nu \approx 0.5 \]
    深度冰山订单与暗藏的做市算法\textbf{像海绵一样吸收冲击},维持系统稳定。
  \end{block}

  \vfill

  \centering\small
  这就是市场冲击的"\textbf{平方根定律}"在微观层面的起源。

\end{frame}

\begin{frame}{VOI——静态存量快照}

  如果说 OFI 是\textbf{流量}(Flow),那 VOI 就是\textbf{存量}(Stock)。

  \[
    \rho_t = \frac{V_t^b - V_t^a}{V_t^b + V_t^a} \in [-1, +1]
  \]

  \vfill

  \begin{columns}[t]
    \column{0.48\textwidth}
    \textcolor{quantgreen}{$\rho_t \to +1$}
    \begin{itemize}
      \item 买单挂单压倒性多于卖单
      \item 价格向上阻力极小
      \item 即将向上跳跃
    \end{itemize}

    \column{0.48\textwidth}
    \textcolor{quantred}{$\rho_t \to -1$}
    \begin{itemize}
      \item 厚重卖盘死死压制
      \item 下方买盘稀薄
      \item 如悬崖落石
    \end{itemize}
  \end{columns}

  \vfill

  \begin{exampleblock}{在最优执行策略中的应用 (Cartea \& Jaimungal 2018)}
    高 VOI 时,最优算法应:\textbf{激进挂出买单}(防踏空)+ \textbf{深藏挂出卖单}(防扫单收割)
  \end{exampleblock}

\end{frame}

\begin{frame}{核心辨析:OFI vs 简单"交易方向"}

  \small
  \begin{tabular}{p{3.0cm}|p{4.2cm}|p{4.5cm}}
    \toprule
    \textbf{对比维度} & \textbf{交易方向 (Trade Direction)} & \textbf{OFI (Order Flow Imbalance)} \\
    \midrule
    数据来源 & 仅已发生的成交数据 & 成交 + 挂单 + \textbf{撤单} \\
    反映对象 & 流动性消耗行为 & 整体供需压迫力 \\
    数据视角 & 静态、被动 & 动态、前瞻 \\
    预测能力 & 有较大噪音 & 强线性关系,稳健 \\
    \bottomrule
  \end{tabular}

  \vfill

  \begin{alertblock}{关键洞察}
    一个极度偏正的 OFI \textbf{不一定意味着巨量市价买入}!它同样可能是因为:
    \begin{itemize}
      \item 空头恐慌性大规模\textbf{撤销卖单}(撤单流贡献)
      \item 多头在最佳买价上\textbf{堆积防御性限价单}
    \end{itemize}
    撤单流是隐藏在传统指标盲区中的\textbf{巨大前瞻性信息}。
  \end{alertblock}

\end{frame}



% ============================================================
\section{总结与实习思考}
% ============================================================

\begin{frame}{核心 Takeaways}

  \begin{enumerate}
    \item \textbf{方向是基础}——没有方向就没有 Alpha;Lee-Ready 与 BVC 各有所长

    \item \textbf{Kyle 模型}——价格冲击 $\lambda$ 是信息风险的均衡定价

    \item \textbf{VPIN 是风控}——体积时钟让我们提前一小时看见闪电崩盘

    \item \textbf{参与者画像}——冰山 / 扫单 / Boehmer 散户构成完整对手盘地图

    \item \textbf{OFI}——线性预测 $\Delta P$ 

    \item \textbf{深度学习}——多级 LOB + CNN/LSTM 突破传统因子的天花板
  \end{enumerate}

  \vfill

  \begin{alertblock}{贯穿全篇的方法论}
    从原始 Tick 数据,层层向上构建结构化的微观特征,最终汇聚成可预测、可执行、可控风险的 Alpha 因子。
  \end{alertblock}

\end{frame}

\begin{frame}{实习落地应用思考}

  这套订单流分类体系如何具体服务于团队的业务?

  \vfill

  \begin{block}{① 算法执行优化}
    \begin{itemize}
      \item 利用 \textbf{VOI} 实时监测盘口压迫力,动态调整子单激进度
      \item 利用\textbf{冰山检测}在跟随机构方向时降低冲击成本
      \item 在高 \textbf{VPIN} 环境下暂停做市,避免逆向选择损失
    \end{itemize}
  \end{block}

  \begin{block}{② 因子挖掘}
    \begin{itemize}
      \item 复用现有 OFI 因子,扩展到多档(Multi-Level OFI)
      \item 探索 \textbf{OFI + 因子中性化}后的稳健性提升
      \item 用 \textbf{BVC} 在缺乏 Tick 数据的市场近似重构毒性指标
    \end{itemize}
  \end{block}

  \begin{block}{③ 风控体系}
    \begin{itemize}
      \item 将 \textbf{VPIN CDF $> 0.9$} 作为系统级流动性预警阈值
      \item 监测撤单/成交比的瞬时跳跃以识别 Spoofing 操纵
    \end{itemize}
  \end{block}

\end{frame}

\begin{frame}{前沿:多级 LOB 与 DeepLOB 深度学习}

  现代高频量化不再满足于单层 L1 失衡——采集 \textbf{LOB 前 10 档甚至 20 档}全深度数据。

  \vfill

  \begin{center}
  \begin{tikzpicture}[node distance=0.6cm, font=\scriptsize]
    \node[draw, rounded corners, fill=blue!10] (a) {Multi-Level LOB Tensor (10+ levels $\times$ OHLC $\times$ Volume)};
    \node[draw, rounded corners, fill=blue!20, below=of a] (b) {1D CNN(提取局部空间特征)};
    \node[draw, rounded corners, fill=blue!30, below=of b] (c) {LSTM(捕捉时序依赖)};
    \node[draw, rounded corners, fill=quantblue!40, below=of c, text=white] (d) {未来 10 个 Tick 的价格反转概率};
    \draw[->] (a) -- (b);
    \draw[->] (b) -- (c);
    \draw[->] (c) -- (d);
  \end{tikzpicture}
  \end{center}

  \vfill

  \textbf{表现}:在预测未来 10 个 Tick 价格反转胜率上,DeepLOB \textbf{远远碾压}传统线性回归统计模型——AI 模型能提取人类难以理解的高维非线性时空相关性。

\end{frame}

\begin{frame}
  \begin{center}
    \Huge 谢谢聆听 \\[1em]
    \Large Q \& A
  \end{center}
\end{frame}

\appendix

\begin{frame}{参考文献(关键文献)}
  \footnotesize
  \begin{itemize}
    \item Lee, C. M. C., \& Ready, M. J. (1991). \emph{Inferring Trade Direction from Intraday Data}. Journal of Finance.
    \item Easley, D., López de Prado, M., \& O'Hara, M. (2012). \emph{Flow Toxicity and Liquidity in a High-Frequency World}. Review of Financial Studies.
    \item Kyle, A. S. (1985). \emph{Continuous Auctions and Insider Trading}. Econometrica.
    \item Easley, D., Kiefer, N. M., O'Hara, M., \& Paperman, J. B. (1996). \emph{Liquidity, Information, and Infrequently Traded Stocks}. Journal of Finance.
    \item Cont, R., Kukanov, A., \& Stoikov, S. (2014). \emph{The Price Impact of Order Book Events}. Journal of Financial Econometrics.
    \item Boehmer, E., Jones, C. M., Zhang, X., \& Zhang, X. (2021). \emph{Tracking Retail Investor Activity}. Journal of Finance.
    \item Cartea, Á., \& Jaimungal, S. (2018). \emph{Algorithmic Trading with Learning}. International Journal of Theoretical and Applied Finance.
    \item Zhang, Z., Zohren, S., \& Roberts, S. (2019). \emph{DeepLOB: Deep Convolutional Neural Networks for Limit Order Books}. IEEE Transactions on Signal Processing.
  \end{itemize}
\end{frame}

\end{document}
